\textbf{Phần 1. Trắc nghiệm lựa chọn}
\setcounterpageref{ex}{0}

\begin{ex} %[5H6V33]
Một mảnh ruộng hình thang có đáy lớn $32\text{ m}$, đáy bé $24\text{ m}$. Nếu mở rộng đáy lớn thêm $4\text{ m}$ thì diện tích mảnh ruộng tăng thêm $36\text{ m}^2$. Diện tích mảnh ruộng ban đầu là
\choice
{$450\text{ m}^2$}
{\True $504\text{ m}^2$}
{$540\text{ m}^2$}
{$252\text{ m}^2$}
\loigiai{
Chiều cao của mảnh ruộng hình thang là:
$36 \times 2 : 4 = 18\text{ (m)}$.\\
Diện tích mảnh ruộng ban đầu là:
$(32 + 24) \times 18 : 2 = 504\text{ (m}^2\text{)}$.\\
\textbf{Đáp án của câu hỏi là B}
}
\end{ex}

\begin{ex} %[5H6V33]
Bánh xe của một xe máy có bán kính $0{,}3\text{ m}$. Khi xe máy đi được quãng đường dài $942\text{ m}$ thì bánh xe đã quay được số vòng là
\choice
{$250\text{ vòng}$}
{\True $500\text{ vòng}$}
{$750\text{ vòng}$}
{$1\,000\text{ vòng}$}
\loigiai{
Chu vi của bánh xe là:
$2 \times 3{,}14 \times 0{,}3 = 1{,}884\text{ (m)}$.\\
Số vòng quay của bánh xe là:
$942 : 1{,}884 = 500\text{ (vòng)}$.\\
\textbf{Đáp án của câu hỏi là B}
}
\end{ex}

\begin{ex} %[5H6V33]
Một hình vuông và một hình tròn có cùng chu vi bằng $62{,}8\text{ cm}$. Khẳng định nào sau đây là đúng?
\choice
{Diện tích hai hình bằng nhau}
{Diện tích hình vuông lớn hơn diện tích hình tròn}
{\True Diện tích hình tròn lớn hơn diện tích hình vuông}
{Không thể so sánh được diện tích}
\loigiai{
Bán kính hình tròn là: $62{,}8 : 3{,}14 : 2 = 10\text{ (cm)}$.\\
Diện tích hình tròn là: $3{,}14 \times 10 \times 10 = 314\text{ (cm}^2\text{)}$.\\
Cạnh hình vuông là: $62{,}8 : 4 = 15{,}7\text{ (cm)}$.\\
Diện tích hình vuông là: $15{,}7 \times 15{,}7 = 246{,}49\text{ (cm}^2\text{)}$.\\
Vì $314 > 246{,}49$ nên diện tích hình tròn lớn hơn diện tích hình vuông.\\
\textbf{Đáp án của câu hỏi là C}
}
\end{ex}

\begin{ex} %[5H6V33]
Một hình tam giác có độ dài đáy $15\text{ cm}$ và chiều cao $12\text{ cm}$. Một hình chữ nhật có diện tích bằng diện tích hình tam giác đó và có chiều rộng là $9\text{ cm}$. Chiều dài của hình chữ nhật là
\choice
{$8\text{ cm}$}
{\True $10\text{ cm}$}
{$12\text{ cm}$}
{$15\text{ cm}$}
\loigiai{
Diện tích hình tam giác là:
$15 \times 12 : 2 = 90\text{ (cm}^2\text{)}$.\\
Chiều dài của hình chữ nhật là:
$90 : 9 = 10\text{ (cm)}$.\\
\textbf{Đáp án của câu hỏi là B}
}
\end{ex}

\begin{ex} %[5H6V33]
Người ta sơn cả hai mặt của một cánh cửa ra vào hình chữ nhật có kích thước $0{,}9\text{ m} \times 2\text{ m}$. Chi phí tiền sơn là $80\,000$ đồng cho mỗi mét vuông. Tổng chi phí sơn cánh cửa đó là
\choice
{$144\,000\text{ đồng}$}
{\True $288\,000\text{ đồng}$}
{$360\,000\text{ đồng}$}
{$200\,000\text{ đồng}$}
\loigiai{
Diện tích một mặt cánh cửa là:
$0{,}9 \times 2 = 1{,}8\text{ (m}^2\text{)}$.\\
Diện tích cần sơn cả hai mặt là:
$1{,}8 \times 2 = 3{,}6\text{ (m}^2\text{)}$.\\
Tổng chi phí tiền sơn là:
$3{,}6 \times 80\,000 = 288\,000\text{ (đồng)}$.\\
\textbf{Đáp án của câu hỏi là B}
}
\end{ex}

\begin{ex} %[5H6V33]
Cho tam giác $ABC$ có cạnh đáy $BC = 20\text{ cm}$. Trên cạnh $BC$ lấy điểm $D$ sao cho $CD = 5\text{ cm}$. Tỉ số diện tích giữa tam giác $ABD$ và tam giác $ABC$ là
\choice
{$\dfrac{1}{4}$}
{$\dfrac{1}{3}$}
{$\dfrac{2}{3}$}
{\True $\dfrac{3}{4}$}
\loigiai{
Độ dài cạnh $BD$ là:
$20 - 5 = 15\text{ (cm)}$.\\
Hai tam giác $ABD$ và $ABC$ có chung đường cao kẻ từ $A$ xuống $BC$ nên tỉ số diện tích bằng tỉ số hai đáy:
$$\dfrac{S_{ABD}}{S_{ABC}} = \dfrac{BD}{BC} = \dfrac{15}{20} = \dfrac{3}{4}.$$
\textbf{Đáp án của câu hỏi là D}
}
\end{ex}

\begin{ex} %[5H6V33]
Một bồn hoa hình tròn nằm trong một cái sân hình vuông cạnh $6\text{ m}$. Biết bồn hoa có đường kính bằng $4\text{ m}$. Diện tích phần sân nằm ngoài bồn hoa là
\choice
{$10{,}88\text{ m}^2$}
{\True $23{,}44\text{ m}^2$}
{$24{,}56\text{ m}^2$}
{$12{,}56\text{ m}^2$}
\loigiai{
Diện tích cái sân hình vuông là:
$6 \times 6 = 36\text{ (m}^2\text{)}$.\\
Bán kính bồn hoa là:
$4 : 2 = 2\text{ (m)}$.\\
Diện tích bồn hoa hình tròn là:
$3{,}14 \times 2 \times 2 = 12{,}56\text{ (m}^2\text{)}$.\\
Diện tích phần sân nằm ngoài bồn hoa là:
$36 - 12{,}56 = 23{,}44\text{ (m}^2\text{)}$.\\
\textbf{Đáp án của câu hỏi là B}
}
\end{ex}

\begin{ex} %[5H6V33]
Một thửa ruộng hình thang có diện tích $450\text{ m}^2$, chiều cao $15\text{ m}$. Biết đáy bé kém đáy lớn $8\text{ m}$. Độ dài đáy bé của thửa ruộng là
\choice
{$22\text{ m}$}
{\True $26\text{ m}$}
{$30\text{ m}$}
{$34\text{ m}$}
\loigiai{
Tổng độ dài hai đáy là:
$450 \times 2 : 15 = 60\text{ (m)}$.\\
Độ dài đáy bé là:
$(60 - 8) : 2 = 26\text{ (m)}$.\\
\textbf{Đáp án của câu hỏi là B}
}
\end{ex}

\begin{ex} %[5H6V33]
Một tấm kính hình thoi có độ dài hai đường chéo là $1{,}2\text{ m}$ và $1{,}5\text{ m}$. Diện tích của tấm kính đó là
\choice
{$1{,}8\text{ m}^2$}
{\True $0{,}9\text{ m}^2$}
{$2{,}7\text{ m}^2$}
{$0{,}45\text{ m}^2$}
\loigiai{
Diện tích tấm kính hình thoi là:
$1{,}2 \times 1{,}5 : 2 = 0{,}9\text{ (m}^2\text{)}$.\\
\textbf{Đáp án của câu hỏi là B}
}
\end{ex}

\begin{ex} %[5H6V33]
Một hình tam giác có độ dài cạnh đáy giảm đi $20\%$ nhưng chiều cao tương ứng tăng thêm $20\%$. Diện tích của hình tam giác đó thay đổi như thế nào?
\choice
{Không thay đổi}
{Tăng thêm $4\%$}
{\True Giảm đi $4\%$}
{Giảm đi $2\%$}
\loigiai{
Gọi đáy là $a$, chiều cao là $h$. Diện tích ban đầu: $S = \dfrac{a \times h}{2}$.\\
Đáy mới: $0{,}8 \times a$; chiều cao mới: $1{,}2 \times h$.\\
Diện tích mới: $S' = \dfrac{(0{,}8 \times a) \times (1{,}2 \times h)}{2} = 0{,}96 \times S = 96\% \times S$.\\
Vậy diện tích giảm đi: $100\% - 96\% = 4\%$.\\
\textbf{Đáp án của câu hỏi là C}
}
\end{ex}

\begin{ex} %[5H6V33]
Bác nông dân rào xung quanh một khu vườn hình chữ nhật dài $25\text{ m}$, rộng $18\text{ m}$. Bác có để một lối vào rộng $2\text{ m}$ không rào. Chiều dài hàng rào bác đã làm là
\choice
{$86\text{ m}$}
{\True $84\text{ m}$}
{$82\text{ m}$}
{$80\text{ m}$}
\loigiai{
Chu vi khu vườn là:
$(25 + 18) \times 2 = 86\text{ (m)}$.\\
Chiều dài hàng rào cần làm là:
$86 - 2 = 84\text{ (m)}$.\\
\textbf{Đáp án của câu hỏi là B}
}
\end{ex}

\begin{ex} %[5H6V33]
Một hình tròn có bán kính $5\text{ cm}$. Nếu tăng bán kính thêm $1\text{ cm}$ thì diện tích hình tròn đó tăng thêm bao nhiêu?
\choice
{$3{,}14\text{ cm}^2$}
{$6{,}28\text{ cm}^2$}
{\True $34{,}54\text{ cm}^2$}
{$31{,}4\text{ cm}^2$}
\loigiai{
Diện tích hình tròn ban đầu: $3{,}14 \times 5 \times 5 = 78{,}5\text{ (cm}^2\text{)}$.\\
Bán kính mới: $5 + 1 = 6\text{ (cm)}$.\\
Diện tích hình tròn mới: $3{,}14 \times 6 \times 6 = 113{,}04\text{ (cm}^2\text{)}$.\\
Diện tích tăng thêm là:
$113{,}04 - 78{,}5 = 34{,}54\text{ (cm}^2\text{)}$.\\
\textbf{Đáp án của câu hỏi là C}
}
\end{ex}

\textbf{Phần 2. Trắc nghiệm đúng sai}
\setcounterpageref{ex}{0}

\begin{ex} %[5H6V33]
Một mảnh đất hình thang có đáy lớn $30\text{ m}$, đáy bé $20\text{ m}$ và chiều cao $12\text{ m}$.
\choiceTF
{\True Diện tích của mảnh đất là $300\text{ m}^2$}
{\True Người ta dùng $30\%$ diện tích mảnh đất để xây nhà, diện tích xây nhà là $90\text{ m}^2$}
{Diện tích phần đất còn lại dùng làm sân và vườn là $200\text{ m}^2$}
{\True Diện tích phần sân và vườn gấp hơn 2 lần diện tích phần xây nhà}
\loigiai{
\begin{itemchoice}
\item Đúng. Diện tích: $(30 + 20) \times 12 : 2 = 300\text{ (m}^2\text{)}$.
\item Đúng. Diện tích làm nhà: $300 \times 30\% = 90\text{ (m}^2\text{)}$.
\item Sai. Diện tích còn lại: $300 - 90 = 210\text{ (m}^2\text{)} 
eq 200\text{ m}^2$.
\item Đúng. Tỉ số là: $210 : 90 \approx 2{,}33 > 2$.
\end{itemchoice}
}
\end{ex}

\begin{ex} %[5H6V33]
Một bánh xe ô tô có đường kính $0{,}8\text{ m}$.
\choiceTF
{\True Bán kính của bánh xe ô tô là $0{,}4\text{ m}$}
{\True Chu vi của bánh xe ô tô là $2{,}512\text{ m}$}
{\True Khi ô tô chạy được quãng đường $251{,}2\text{ m}$ thì bánh xe lăn đúng $100$ vòng}
{Khi xe chạy với vận tốc $50{,}24\text{ km/h}$ trong 1 giờ thì bánh xe lăn được đúng $2\,000$ vòng}
\loigiai{
\begin{itemchoice}
\item Đúng. Bán kính: $0{,}8 : 2 = 0{,}4\text{ (m)}$.
\item Đúng. Chu vi: $0{,}8 \times 3{,}14 = 2{,}512\text{ (m)}$.
\item Đúng. Số vòng quay: $251{,}2 : 2{,}512 = 100\text{ (vòng)}$.
\item Sai. Quãng đường đi được trong 1 giờ là $50{,}24\text{ km} = 50\,240\text{ m}$. Số vòng quay là: $50\,240 : 2{,}512 = 20\,000\text{ (vòng)} 
eq 2\,000\text{ vòng}$.
\end{itemchoice}
}
\end{ex}

\textbf{Phần 3. Trắc nghiệm trả lời ngắn}
\setcounterpageref{ex}{0}

\begin{ex} %[5H6V33]
Một hình tam giác có diện tích bằng $135\text{ cm}^2$ và chiều cao là $15\text{ cm}$. Độ dài cạnh đáy của hình tam giác đó là bao nhiêu xăng-ti-mét?
\shortans{18}
\loigiai{
Độ dài cạnh đáy hình tam giác là:
$135 \times 2 : 15 = 18\text{ (cm)}$.
}
\end{ex}

\begin{ex} %[5H6V33]
Một hình thang có diện tích bằng $144\text{ m}^2$ và chiều cao là $8\text{ m}$. Trung bình cộng độ dài hai đáy của hình thang đó bằng bao nhiêu mét?
\shortans{18}
\loigiai{
Trung bình cộng hai đáy của hình thang là:
$144 : 8 = 18\text{ (m)}$.
}
\end{ex}

\begin{ex} %[5H6V33]
Một bánh xe đạp có đường kính $0{,}7\text{ m}$. Hỏi bánh xe phải quay bao nhiêu vòng để xe đi được quãng đường dài $2\,198\text{ m}$?
\shortans{1000}
\loigiai{
Chu vi bánh xe là:
$0{,}7 \times 3{,}14 = 2{,}198\text{ (m)}$.\\
Số vòng quay cần thiết là:
$2\,198 : 2{,}198 = 1\,000\text{ (vòng)}$.
}
\end{ex}

\begin{ex} %[5H6V33]
Một khu vườn hình chữ nhật có chu vi $120\text{ m}$. Chiều dài gấp đôi chiều rộng. Diện tích của khu vườn đó là bao nhiêu mét vuông?
\shortans{800}
\loigiai{
Nửa chu vi khu vườn là:
$120 : 2 = 60\text{ (m)}$.\\
Chiều rộng khu vườn là:
$60 : (1 + 2) = 20\text{ (m)}$.\\
Chiều dài khu vườn là:
$20 \times 2 = 40\text{ (m)}$.\\
Diện tích khu vườn là:
$40 \times 20 = 800\text{ (m}^2\text{)}$.
}
\end{ex}

\begin{ex} %[5H6V33]
Một hình thang vuông có chiều cao bằng $10\text{ cm}$, độ dài hai đáy lần lượt là $8\text{ cm}$ và $16\text{ cm}$. Diện tích của hình thang vuông đó là bao nhiêu xăng-ti-mét vuông?
\shortans{120}
\loigiai{
Diện tích hình thang vuông là:
$(8 + 16) \times 10 : 2 = 120\text{ (cm}^2\text{)}$.
}
\end{ex}

\begin{ex} %[5H6V33]
Một tấm bìa hình vuông có cạnh dài $20\text{ cm}$. Tính diện tích của tấm bìa hình vuông đó theo đơn vị xăng-ti-mét vuông.
\shortans{400}
\loigiai{
Diện tích tấm bìa hình vuông là:
$20 \times 20 = 400\text{ (cm}^2\text{)}$.
}
\end{ex}
